\chapter{Introduction}

\section{Why the Lambek calculus matters}
The Lambek calculus was introduced in \citet{Lambek1958} as a formal calculus of syntactic
types. Its original purpose was already recognisably linguistic: to give a principled method for
distinguishing well-formed from ill-formed strings and to model the way in which syntactic
categories combine in natural language. Lambek's follow-up work in \citet{Lambek1961}
further developed the type-theoretic viewpoint and helped establish categorial grammar as a
serious mathematical approach to syntax.

The importance of the Lambek calculus lies in the fact that it sits at an unusually fruitful
intersection of logic, algebra, and linguistics. On the logical side, it is one of the prototypical
substructural logics: order matters, and structural principles such as exchange, weakening, and
contraction are not built in by default. On the linguistic side, this resource sensitivity is not an
artificial restriction but an advantage, because word order and local composition are central to
syntactic analysis. For this reason the Lambek calculus became a foundational system in the
study of categorial and type-logical grammars \citep{vanBenthem1988,Moortgat1995Multimodal,morrill1996generalising}.

A second reason for its importance is that the calculus offers an explicit bridge between syntax
and semantics. Derivations in categorial calculi naturally correspond to compositional analyses of
meaning, and this connection has been one of the main reasons for the enduring influence of
Lambek-style systems in formal semantics and mathematical linguistics \citep{vanBenthem1988,Steedman2000}.


\section{Completeness: classical results and the open problem of this thesis}
There is an important distinction to make at the outset. The completeness of the \emph{classical}
Lambek calculus is not an open problem in general. Several major completeness results are
already known. \citet{Dosen1985} proved a completeness theorem for the Lambek calculus of
syntactic categories; \citet{AndrekaMikulas1994} established strong completeness results for
relational semantics; and \citet{Pentus1995} proved completeness with respect to language
models. In addition, Pentus' work on the context-free status of Lambek grammars clarified the
place of the calculus in formal-language theory \citep{Pentus1993,Pentus1997}.

The open problem relevant to the present thesis is therefore narrower and more specific:
\emph{is the probabilistic semantics introduced here complete for the Lambek calculus?} This
question is open because the semantics developed in this thesis is corpus-relative, numerical,
and intentionally partial in the sense that it is tied to observed linguistic data. The standard
canonical-model and algebraic techniques from classical completeness proofs do not transfer
immediately to such a setting.

For that reason the present thesis focuses on soundness first. Establishing soundness is the right
first milestone: before asking whether the semantics captures \emph{all} valid inferences, one
must first show that every derivable Lambek inequality is semantically valid. Completeness is the
next goal, not a claim of the present work.


\section{Lambek calculus in linguistics}
In linguistics, the Lambek calculus is best understood as part of the broader tradition of
categorial grammar. Categories behave like types, and grammatical composition is governed by
logical operations rather than by phrase-structure rules. This makes the calculus attractive for
analysing phenomena where lexical information and local combinatorics are central.

Type-logical grammar developed this perspective in depth and connected it to proof theory,
formal semantics, and computational linguistics \citep{Moortgat1997,Morrill1994}. Closely related
systems, especially Combinatory Categorial Grammar (CCG), have been particularly successful in
wide-coverage parsing and semantic interpretation \citep{Steedman2000,ClarkCurran2007}.
Although CCG is not identical to the Lambek calculus, the family resemblance is close enough
that probabilistic work on CCG provides a useful source of intuition for the present project.


\section{Motivating Probabilistic Semantics}
Pure derivability is not enough for realistic parsing. A sentence may have several syntactically
well-formed derivations, and those derivations may correspond to different readings. This is not a
pathology but a normal feature of natural language. A parser that returns every derivation with
equal status is therefore mathematically respectable but practically incomplete.

Probabilistic grammatical models address this by associating weights or probabilities with lexical
choices and structural combinations. In probabilistic context-free grammars, probabilities rank
parse trees generated by grammar rules; in probabilistic CCG, similar ideas are applied to a more
lexicalised grammar formalism \citep{ManningSchutze1999,ClarkCurran2007}. There has also been
explicit prior work on stochastic Lambek grammars \citep{BonfanteDeGroote2004}. The basic
intuition is simple: if several analyses are licensed, prefer the one better supported by the data.

This thesis adopts that intuition in a deliberately conservative way. Rather than starting from a
full statistical parser, it starts from a semantic valuation $v(x,A) \in [0,1]$ for strings $x$ and
Lambek types $A$, interpreted relative to a corpus. The main question is then whether the
algebraic behaviour of the Lambek connectives can be captured in a way that supports a
sound proof theory.


\section{Aims and scope of the thesis}
The thesis has four main aims.
\begin{enumerate}[label=(\roman*)]
    \item To motivate a probabilistic semantics for the associative Lambek calculus from the
    perspective of parsing and ambiguity resolution.
    \item To formalise that semantics carefully enough that the main recursive clauses are
    mathematically well-defined.
    \item To prove soundness of the Lambek axioms with respect to the proposed semantics.
    \item To identify the obstacles that make completeness the next major step rather than a result
    already achieved.
\end{enumerate}

Several things are \emph{not} attempted here. The thesis does not provide a learning algorithm for
extracting lexical probabilities from raw corpora, does not present a full parser implementation,
and does not prove completeness. These tasks are deferred to future work.

